% Section 5: Discussion (1,200 words target)
% TODO: Complete after results interpretation

\subsection{Interpretation of Findings}

% If findings are robust (H0):
% - Strengthens portfolio claims
% - Validates tract-level as appropriate resolution
% - Implications for practitioners

% If findings are sensitive (H1):
% - MAUP is real concern
% - Need resolution-aware analysis
% - Implications for methodology standards

\subsection{Theoretical Implications}

\paragraph{MAUP Mechanisms.}
% TODO: Explain why resolution matters (or doesn't)
% - Geographic clustering at different scales
% - Minority community cohesion vs dispersion
% - Boundary granularity effects on compactness

\paragraph{Algorithm Behavior.}
% TODO: Discuss METIS interaction with graph size
% - Edge-cut minimization at different scales
% - Local optima and solution diversity

\subsection{Policy Implications}

\paragraph{For Redistricting Commissions.}
% TODO: Guidance on resolution choice
% If robust → use tracts (efficiency)
% If sensitive → use finest feasible resolution

\paragraph{For Courts.}
% TODO: VRA compliance implications
% Resolution specification in legal standards
% Majority-minority definition depends on scale

\subsection{Limitations}

\begin{itemize}
    \item \textbf{Computational}: Block-level analysis may face memory constraints for largest states
    \item \textbf{Data quality}: Blocks have small sample sizes, introducing demographic noise
    \item \textbf{Scope}: Only recursive bisection tested; other partitioners may behave differently
    \item \textbf{Census years}: Focus on 2020; temporal consistency untested
\end{itemize}

\subsection{Future Work}

Promising directions include:
\begin{itemize}
    \item Testing MAUP with other partitioning algorithms (SCOTCH, spectral methods)
    \item Examining different adjacency structures (queen contiguity, distance-based)
    \item Multi-resolution ensemble approaches
    \item Extension to multi-census temporal analysis
\end{itemize}
